\documentclass[11pt,a4paper]{article}
\usepackage[utf8]{inputenc}
\usepackage[T1]{fontenc}
\usepackage[french]{babel}
\usepackage[margin=2.2cm]{geometry}
\usepackage{xcolor}
\usepackage{listings}
\usepackage{enumitem}
\usepackage{titlesec}
\usepackage{hyperref}
\hypersetup{colorlinks=true, linkcolor=blue, urlcolor=blue}

\definecolor{bleu}{HTML}{1d6ff2}
\definecolor{rouge}{HTML}{e03131}
\definecolor{gris}{HTML}{f0f0f0}

\lstset{basicstyle=\ttfamily\footnotesize, backgroundcolor=\color{gris}, frame=single, breaklines=true, language=PHP, showstringspaces=false}

\titleformat{\section}{\Large\bfseries\color{bleu}}{\thesection}{1em}{}
\titleformat{\subsection}{\large\bfseries\color{rouge}}{\thesubsection}{1em}{}

\newcommand{\question}[1]{\vspace{0.4em}\noindent\textbf{\color{rouge}Q : #1}\par}
\newcommand{\reponse}[1]{\noindent\textbf{R :} #1\par\vspace{0.4em}}
\newcommand{\motcle}[1]{\noindent\textbf{\color{bleu}#1} :\ }

\title{\textbf{Guide complet de soutenance -- LO07 BlaBlaCar 2026}\\[0.3em]\large Partie de Pierre : la couche \textbf{Vue} et le visuel (Bootstrap)}
\author{Pierre Bheidi \quad (binôme : Moss'Ab Mirande-Ney)}
\date{}

\begin{document}
\maketitle
\tableofcontents
\vspace{1em}
\hrule
\vspace{1em}

\noindent\textit{Ce document est un guide complet : même sans avoir codé cette
partie, tu pourras tout présenter et répondre aux questions du jury. Lis-le une
fois en entier, puis relis le « script de démo » et les « questions/réponses ».}

\section{Contexte du projet (à connaître absolument)}

\motcle{Le sujet} BlaBlaCar est une application web de \textbf{covoiturage}. Les
utilisateurs se connectent et, selon leur \textbf{rôle}, accèdent à des
fonctionnalités différentes :
\begin{itemize}[noitemsep]
  \item \textbf{Administrateur} : gère les utilisateurs, les véhicules, les villes.
  \item \textbf{Conducteur} : gère ses véhicules et ses trajets (création, clôture).
  \item \textbf{Passager} : réserve des trajets et consulte ses réservations.
\end{itemize}

\motcle{La base de données} 5 tables : \texttt{utilisateur} (id, nom, prenom,
role, login, password, solde), \texttt{vehicule}, \texttt{ville},
\texttt{trajet} (départ, arrivée, conducteur, véhicule, prix, date, statut
actif/passif), \texttt{reservation} (lie un trajet et un passager).

\motcle{L'architecture} Le projet suit le patron \textbf{MVC}
(Modèle-Vue-Contrôleur) : on sépare les données (Modèle), l'affichage (Vue) et la
logique (Contrôleur). \textbf{Ma partie, c'est la Vue.}

\section{Les concepts que je dois savoir expliquer}

\motcle{MVC} Trois couches. Le \textbf{Modèle} parle à la base. Le
\textbf{Contrôleur} reçoit la requête, appelle le modèle, prépare les données. La
\textbf{Vue} (ma partie) affiche ces données en HTML. Avantage : on peut changer
l'affichage sans toucher à la logique.

\motcle{Une « vue »} Un fichier PHP qui contient surtout du HTML, avec quelques
petites boucles pour afficher une liste. Une vue \textbf{ne fait jamais de requête
SQL} ni de calcul métier : elle reçoit des données déjà prêtes.

\motcle{Bootstrap} Une bibliothèque CSS/JS gratuite qui fournit des composants
déjà stylés (barre de navigation, cartes, tableaux, boutons). On l'inclut par un
lien CDN ; on applique ensuite des \textbf{classes} sur nos balises HTML
(ex.~\texttt{class="btn btn-primary"}). On a pris la version \textbf{5.3}, la même
que dans les TD.

\motcle{Session} Une mémoire côté serveur qui retient qui est connecté entre les
pages (variable \texttt{\$\_SESSION['login\_id']}). La Vue la lit pour savoir quel
menu afficher.

\motcle{Échappement (htmlspecialchars)} Transformer les caractères spéciaux HTML
(\texttt{<}, \texttt{>}) en texte inoffensif avant l'affichage, pour empêcher
l'injection de code (faille \textbf{XSS}).

\section{Ma partie en détail : les fichiers de la Vue}

Tous mes fichiers sont dans \texttt{projet/app/views/}. Ils sont rangés par rôle :
\texttt{admin/}, \texttt{conducteur/}, \texttt{passager/}, \texttt{examinateur/},
\texttt{innovations/}, \texttt{auth/}, \texttt{home/}, et un dossier
\texttt{layout/} commun.

\subsection{Le gabarit commun (layout/)}
Trois fragments composent l'ossature de \textbf{chaque} page :
\begin{itemize}[noitemsep]
  \item \texttt{header.php} : le début du HTML (\texttt{<head>}), le chargement de
        Bootstrap (CSS + JS par CDN) et un petit bloc de style (le conteneur a une
        marge haute pour passer sous la barre fixe).
  \item \texttt{fragmentMenu.php} : la \textbf{barre de navigation} et le bandeau
        bleu « Projet BlaBlaCar 2026 » (le \emph{jumbotron}).
  \item \texttt{footer.php} : le pied de page (noms, UTT, date) et la fermeture du HTML.
\end{itemize}
C'est le contrôleur qui les assemble via une méthode \texttt{render()} : on n'écrit
donc jamais deux fois l'en-tête.

\subsection{La barre de menu dynamique (le point fort)}
C'est la pièce centrale de ma partie. Elle :
\begin{enumerate}[noitemsep]
  \item lit l'utilisateur connecté en session ;
  \item affiche en haut à gauche les \textbf{noms du binôme}, et à droite le
        \textbf{nom de l'utilisateur connecté} avec son \textbf{solde} ;
  \item affiche \textbf{un seul} menu métier selon le rôle (Administrateur,
        Conducteur \emph{ou} Passager) ;
  \item affiche toujours les menus Innovations, Examinateur et Se connecter.
\end{enumerate}
\begin{lstlisting}
if ($courant !== null && $courant['role'] === 'administrateur'): ...
elseif ($courant['role'] === 'conducteur'): ...
elseif ($courant['role'] === 'passager'): ...
endif;
\end{lstlisting}
« Si personne n'est connecté, aucun menu métier n'apparaît » -- exactement ce que
demande le sujet.

\subsection{Le style Bootstrap utilisé (identique aux TD du prof)}
\begin{itemize}[noitemsep]
  \item \textbf{Barre} : \texttt{navbar navbar-expand-lg bg-dark navbar-dark
        fixed-top} avec des menus déroulants (\texttt{dropdown}).
  \item \textbf{Bandeau} : \texttt{<div class="mt-4 p-5 bg-primary text-white
        rounded">} (le jumbotron moderne de Bootstrap 5).
  \item \textbf{Contenu} : chaque page est une \textbf{carte}
        (\texttt{card / card-body bg-info}) avec un titre, et le contenu dans un
        bloc intérieur \texttt{bg-light p-3 rounded}.
  \item \textbf{Tableaux} : \texttt{table table-striped table-bordered
        table-hover} avec un en-tête rouge \texttt{thead table-danger}.
  \item \textbf{Formulaires} : \texttt{form-control}, \texttt{form-select},
        boutons \texttt{btn btn-primary}.
  \item \textbf{Messages} : \texttt{alert alert-success} (succès) ou
        \texttt{alert-danger} (erreur) ; cartes \texttt{text-bg-success} /
        \texttt{text-bg-danger} pour les résultats d'opération.
\end{itemize}
Tout est en Bootstrap, \textbf{sans fichier CSS personnel} (juste un petit bloc de
style inline pour le conteneur, comme dans les TD).

\subsection{Quelques vues concrètes}
\begin{itemize}[noitemsep]
  \item \texttt{home/accueil.php} : page d'accueil + comptes de test.
  \item \texttt{auth/login.php} : formulaire de connexion.
  \item \texttt{admin/utilisateurs.php} : tableau de tous les utilisateurs.
  \item \texttt{admin/vehicules.php} : tableau des véhicules. \textbf{Attention} :
        on n'affiche \emph{pas} les clés primaires, et le propriétaire est
        reconstruit « prénom + nom » -- une exigence du sujet.
  \item \texttt{conducteur/trajets.php} : trajets avec un \textbf{badge} vert
        (actif) ou gris (passif).
  \item \texttt{innovations/data.php} : le tableau de bord (statistiques).
\end{itemize}

\section{Déroulé : que se passe-t-il quand une page s'affiche ?}
\begin{enumerate}[noitemsep]
  \item L'utilisateur clique sur un lien, ex. \texttt{index.php?action=admin\_utilisateurs}.
  \item Le contrôleur récupère les données (la liste des utilisateurs) auprès du modèle.
  \item Il appelle \texttt{render('admin/utilisateurs', \$data)}.
  \item \texttt{render()} assemble header + fragmentMenu + ma vue + footer, et
        \texttt{extract(\$data)} rend la variable \texttt{\$utilisateurs}
        disponible dans la vue.
  \item Ma vue parcourt \texttt{\$utilisateurs} avec un \texttt{foreach} et génère
        les lignes du tableau, chaque valeur passée par \texttt{htmlspecialchars}.
\end{enumerate}

\section{Script de démo (ce que je dis à l'écran)}
« Voici la page d'accueil. En haut, la barre de navigation affiche nos deux noms.
Tant qu'on n'est pas connecté, seuls trois menus apparaissent : Innovations,
Examinateur et Se connecter. \emph{[Je me connecte.]} Maintenant qu'on est
administrateur, un nouveau menu Administrateur est apparu, et à droite s'affichent
le nom et le solde. \emph{[Liste des utilisateurs.]} C'est une carte Bootstrap avec
un tableau à en-tête rouge. \emph{[Liste des véhicules.]} Remarquez qu'on n'affiche
pas les identifiants techniques : le propriétaire est reconstruit à partir du
prénom et du nom. \emph{[Je réduis la fenêtre.]} La barre se replie en menu
hamburger : c'est le responsive de Bootstrap. »

\section{Questions potentielles du jury (et réponses)}

\question{C'est quoi la couche Vue ?}
\reponse{La partie qui s'occupe uniquement de l'affichage. Elle reçoit des données
déjà préparées par le contrôleur et les met en forme en HTML, sans logique métier
ni requête à la base.}

\question{Comment la barre de menu sait quel menu afficher ?}
\reponse{Elle lit \texttt{\$\_SESSION['login\_id']}, récupère l'utilisateur via le
modèle, puis teste son champ \texttt{role}. Un seul des trois menus métier est
affiché selon ce rôle.}

\question{Et si personne n'est connecté ?}
\reponse{\texttt{login\_id} vaut -1, donc aucun menu métier n'est affiché : seuls
Innovations, Examinateur et Se connecter restent visibles.}

\question{Pourquoi Bootstrap ?}
\reponse{Pour un rendu propre et responsive immédiat, avec des composants testés.
C'est aussi la bibliothèque utilisée dans les TD. On garde la maîtrise du HTML et
des classes choisies.}

\question{Avez-vous écrit du CSS ?}
\reponse{Quasiment pas : on utilise les classes Bootstrap. Il y a seulement un petit
bloc de style inline dans l'en-tête pour la marge du conteneur sous la barre fixe.}

\question{À quoi sert \texttt{htmlspecialchars} ?}
\reponse{À échapper les caractères HTML en sortie : c'est une protection contre les
failles XSS. Du code injecté dans un champ s'affiche comme du texte au lieu d'être
exécuté.}

\question{Comment une vue reçoit-elle ses données ?}
\reponse{Le contrôleur passe un tableau associatif à \texttt{render()}, qui fait un
\texttt{extract()} : chaque clé devient une variable utilisable dans la vue.}

\question{Différence entre \texttt{<?= } et \texttt{<?php echo} ?}
\reponse{\texttt{<?= \$x ?>} est simplement un raccourci d'écriture pour
\texttt{echo \$x}. On s'en sert pour afficher une valeur dans le HTML.}

\question{C'est quoi un fragment / le fragmentMenu ?}
\reponse{Un morceau de vue réutilisé sur toutes les pages. Le fragmentMenu est la
barre de navigation ; on l'écrit une seule fois et il est inclus partout par
\texttt{render()}.}

\question{Comment gérez-vous l'affichage mobile ?}
\reponse{Avec les classes responsive de Bootstrap : \texttt{navbar-toggler} replie
le menu en hamburger, et les tableaux peuvent défiler horizontalement.}

\question{Qui gère l'ouverture des menus déroulants ?}
\reponse{Le JavaScript de Bootstrap (le bundle chargé dans l'en-tête). On ajoute
juste les attributs \texttt{data-bs-toggle="dropdown"}.}

\question{Pourquoi reconstruire le propriétaire « prénom + nom » dans la liste des
véhicules ?}
\reponse{Parce que le sujet interdit d'afficher les clés primaires : on ne montre
pas l'identifiant du propriétaire mais son nom lisible, obtenu par une jointure
faite dans le modèle.}

\question{Y a-t-il de la logique métier dans vos vues ?}
\reponse{Non. Les vues ne contiennent que des boucles et conditions d'affichage.
Tout le métier est dans les contrôleurs et les modèles.}

\end{document}
